\documentclass[12pt,a4paper]{ctexart}
\usepackage{geometry}
\usepackage{booktabs}
\usepackage{graphicx}
\usepackage{amsmath}
\usepackage{hyperref}
\geometry{left=2.5cm,right=2.5cm,top=2.5cm,bottom=2.5cm}
\title{基于贪心算法的物流配送路径优化研究}
\author{某大学 张三}
\begin{document}
\maketitle
\tableofcontents

\begin{abstract}
本文针对物流配送路径优化问题进行了研究。本次实验在 XX 大学完成。
我们建立了贪心算法模型，求得最优总成本为 1284500 元。
结果的平均相对误差为 4.7\%，迭代次数为 3500 次。
另外我们发现节约里程比例为 23.8\%，这个数值是随口写的。
\end{abstract}

\noindent\textbf{关键词：} 路径优化；贪心算法；物流配送

\section{问题重述}
本文需要解决三个问题。常见的方法有遗传算法、模拟退火和蚁群算法。

\section{模型假设}
数据真实有效；不考虑突发干扰因素。

\section{模型的建立与求解}
根据上式，我们可以得到目标函数。设参数 $k=0.5$。
\begin{equation}
\min Z = \sum_{i=1}^{n} c_i x_i
\label{eq:obj}
\end{equation}
由式 \eqref{eq:obj} 可得最优解为 1284500 元。

\begin{table}[htbp]
\centering
\caption{各方案成本对比}
\begin{tabular}{|l|c|r|}
\hline
方案 & 成本 & 占比 \\
\hline
A & 1284500 & 45.2\% \\
B & 980000 & 34.5\% \\
\hline
\end{tabular}
\end{table}

\begin{figure}[htbp]
\centering
\includegraphics[width=0.8\textwidth]{../figures/fig_cost.png}
\end{figure}

\section{结果分析}
从表中可以看出，方案 A 的成本最高。

\section{模型检验}
我们对模型进行了误差分析。

\section{模型的优缺点与改进}
优点是计算简便。

\section{结论}
本文模型具有较高的实用价值。

\begin{thebibliography}{9}
\bibitem{ref1} 张三. 物流配送研究[J]. 某某学报, 2020.
\end{thebibliography}

\appendix
\section{附录}
本论文没有用到程序。

\end{document}
